\documentclass[12pt,a4paper,oneside]{book}
\usepackage[utf8]{inputenc}
\usepackage[T1]{fontenc}
\usepackage[english,french]{babel}
\usepackage{geometry}
\usepackage[table]{xcolor}

\geometry{top=2.5cm,bottom=2.5cm,left=2.5cm,right=2.5cm}
\definecolor{bleuIA}{RGB}{37, 99, 235}

\begin{document}

\begin{titlepage}
\begin{center}
\vspace*{3cm}
{\Huge\bfseries\color{bleuIA}E-GRAPHISME}\\[0.3cm]
{\LARGE by ELECTRON}\\[2cm]
{\huge AI AUTOMATION CENTER}\\[0.3cm]
{\huge ORCHESTRATEUR CENTRAL DES AGENTS IA}\\[1cm]
{\Large Volume 15}\\[3cm]
{\Large Version 2026}
\end{center}
\end{titlepage}

\chapter{Vision}
L'automatisation ne remplace pas les utilisateurs.

\chapter{Agents IA}
\begin{itemize}
\item Commercial AI
\item Marketing AI
\item Finance AI
\item RH AI
\item Support AI
\item Designer AI
\item Developer AI
\end{itemize}

\chapter{AI Orchestrator}
Réception des demandes, identification des agents, répartition des tâches.

\chapter{Gestion des Workflows}
Création de workflows graphiques.

\chapter{Planificateur}
Exécutions immédiates, tâches programmées.

\chapter{Quality Manager}
Contrôle de conformité.

\chapter{AI Mission Control}
Centre de mission pour objectifs complexes.

\chapter{AI Knowledge Router}
Routeur de connaissances.

\chapter{Tableau de Bord d'Orchestration}
Supervision des agents IA.

\end{document}
